% -*- coding: utf-8

\documentclass[slideopt,A4]{beamer}

\usetheme{default}

\useoutertheme{infolines}

\usepackage[]{fontspec}
\usepackage{fontenc}
\usepackage[type1]{libertine}
\renewcommand{\ttdefault}{cmtt}

\usepackage{lastpage}
\usepackage{mdframed}
\usepackage[usenames,dvipsnames,svgnames,table]{xcolor}


\usepackage{eurosym}
\usepackage{url}
\usepackage{anyfontsize}
\setmainfont{linuxbiolinumo}
\setsansfont{linuxbiolinumo}

\definecolor{fontmaincolor}{HTML}{000000}
\definecolor{bgcolor}{HTML}{FFFFFF}
\definecolor{bgshadecolor}{HTML}{DDDDDD}
\graphicspath{%
  {./Figures/}%
}
%\usepackage [T1] {inputenc}
%\usepackage [utf8] {inputenc}

\usepackage{tikz}
\usetikzlibrary{arrows,patterns,plotmarks,shapes,snakes,er,3d,automata,backgrounds,topaths,trees,petri,mindmap}
\usepackage{pgfplots}
\usepackage{pgfplotstable}

\usepackage{eso-pic}
\usepackage{../monster2e}
\usepackage{tikz}
\usepackage{layout}
\usepackage{xcolor}
\usepackage{amsmath,amssymb,wasysym}
\usepackage{fancybox}
\usepackage{graphicx}
\usepackage{multimedia}
\usepackage[absolute,showboxes,overlay]{textpos}
\TPshowboxesfalse
\textblockorigin{0mm}{0mm}

\usepackage{media9}
%\usepackage{movie15_dvipdfmx}
%\usepackage[dvipdfmx]{movie15_dvipdfmx}

\definecolor{rouge1}{RGB}{226,0,38}  % red P
\definecolor{orange1}{RGB}{243,154,38}  % orange P
\definecolor{jaune}{RGB}{254,205,27}  % jaune P
\definecolor{blanc}{RGB}{255,255,255} % blanc P

\definecolor{rouge2}{RGB}{230,68,57}  % red S
\definecolor{orange2}{RGB}{236,117,40}  % orange S
\definecolor{taupe}{RGB}{134,113,127} % taupe S
\definecolor{gris}{RGB}{91,94,111} % gris S
\definecolor{bleu1}{RGB}{38,109,131} % bleu S
\definecolor{bleu2}{RGB}{28,70,114} % bleu S
\definecolor{vert1}{RGB}{133,146,66} % vert S
\definecolor{vert2}{RGB}{157,193,7} % vert S

\definecolor{rouge3}{RGB}{255,200,130}  

\setbeamertemplate{navigation symbols}{}
\setbeamertemplate{blocks}[rounded][shadow=false]
\setbeamercolor{block title}{fg=white, bg=rouge1}
\setbeamercolor{block body}{fg=black, bg=rouge3}

\usepackage{tikz}
% tikz library (frise lncmi exemple)
\usetikzlibrary{calc,fit,shadows,arrows,positioning}
\usepackage{xcolor}
\xdefinecolor{fColor}{RGB}{0 81 145} % Light blue
\newcommand{\fColor}{fColor}

\newcommand{\topLegend}[2]{
  node [text centered,%
  shade,%
  top color=\fColor,%
  text justified,
  draw,%
  semithick,
  rounded corners=2pt,
  above,
  #1] {#2}}

\newcommand{\bottomLegend}[2]{
  node [text centered,%
  shade,%
  top color=\fColor,%
  text justified,
  draw,%
  semithick,
  rounded corners=2pt,
  below,
  #1] {#2}}

\newcommand{\topEvent}[3]{\draw[<-,very thick] (#1) coordinate (#2) (#2 |- xH axis) -- ++ (#3)}
\newcommand{\bottomEvent}[3]{\draw[<-,very thick] (#1) coordinate (#2) -- ++ (#3)}
\newcommand{\bottomGraphics}[2]{node[below,draw,inner sep=1pt,thin,rounded corners=1pt] {\includegraphics[#1]{#2}}}
\newcommand{\topGraphics}[2]{node[above,draw,inner sep=1pt,thin,rounded corners=1pt] {\includegraphics[#1]{#2}}}

% define trademark symbol
\def\registered{{\ooalign{\hfil\raise .00ex\hbox{\scriptsize R}\hfil\crcr\mathhexbox20D}}}

\renewcommand{\sfdefault}{lmss}
\sffamily

\setbeamersize{text margin left=1cm,text margin right=1cm}


\AtBeginSection[] {
  \begin{frame}<beamer>{Plan}
   \tableofcontents[sectionstyle=show/shaded,subsectionstyle=show/shaded/hide]
%    \tableofcontents[currentsection]
  \end{frame}
 }

\AtBeginSubsection[] {
  \begin{frame}<beamer>{Plan}
   \tableofcontents[sectionstyle=show/shaded,subsectionstyle=show/shaded/hide]
%    \tableofcontents[currentsection,currentssubection]
  \end{frame}
 }


\newcommand{\ds}{\displaystyle}
\newcommand{\dg}{^\circ}
\newcommand{\saut}{\vspace*{3mm}\\}
\newcommand{\RR}{\hbox{\bf I\hspace*{-1mm}R}}
%% Macros
\newcommand{\ddx}[1]{\ensuremath{\frac{\partial #1}{\partial x}}}
\newcommand{\ddy}[1]{\ensuremath{\frac{\partial #1}{\partial y}}}
\newcommand{\ddz}[1]{\ensuremath{\frac{\partial #1}{\partial z}}}
\newcommand{\dddxdx}[1]{\ensuremath{\frac{\partial^2 #1}{\partial x \partial x}}}
\newcommand{\dddxdy}[1]{\ensuremath{\frac{\partial^2 #1}{\partial x \partial y}}}
\newcommand{\dddydy}[1]{\ensuremath{\frac{\partial^2 #1}{\partial y \partial y}}}
\renewcommand{\div}{\operatorname{div}}
\renewcommand{\rot}{\operatorname{rot}}

\begin{document}


\setbeamertemplate{background canvas}[vertical shading][top=white,middle=white,bottom=white]

\setbeamertemplate{footline}{ \hspace{5em} \textcolor{white} {Les Mathématiques}\hspace{2em}\null \vspace*{3pt}}


\begin{frame}

\begin{textblock*}{15cm}(13mm,30mm)
{\textcolor{red} {
{\huge\bf Master 2 Recherche: }\\[2mm]
{\bf Équations aux dérivées partielle}\\[8mm] }}
{\textcolor{black} {
%   	{\Large Philippe Helluy \& Christophe Prud'homme}\\[2mm]
	{\Large Université de Strasbourg}\\[2mm]
        \includegraphics[height=1.2cm]{LOGOS/logoCemosis}	\hspace*{40mm}	
        \includegraphics[height=1.2cm]{LOGOS/logoUDS}\\
      }}
\end{textblock*}

\end{frame}


%%%%%%%%%%%%%%%%%%%%%%%%%%%

\setbeamertemplate{background canvas}[vertical shading][top=blanc, middle=blanc,bottom=blanc]

\setbeamercolor{toto}{fg=white,bg=blue}

\setbeamertemplate{footline}
{
\begin{beamercolorbox}[wd=1\paperwidth,ht=15.5pt]{toto}
\hspace{-1.6mm}	
  \raisebox{1.2ex}
  {  \includegraphics[height=.6cm]{LOGOS/logoCemosis}}
  \raisebox{2.5ex}
  %{ P. Helluy \& C. Prud'homme (UNISTRA/CEMOSIS)  }
  {Cemosis, IRMA, UFR Math-Info}
\hspace{\fill}	
 % \raisebox{1.2ex}
%\hspace{\fill}	
%  \raisebox{2.5ex}
%  { \insertframenumber \hspace{1mm}  }
  \raisebox{2.5ex}
 {  }
  \raisebox{1ex}
{  \colorbox{white}{\includegraphics[height=.6cm]{LOGOS/logoUDS}}}
\end{beamercolorbox} 
}

%
%
%%%%%%%%%%%%%%%%%%%%%%%%%%%%%%%%%%%%%%

\begin{frame}{{\Large Plan}}
  \tableofcontents
  % Vous pouvez, si vous le souhaiter ajouter l'option [pausesections]
\end{frame}



\end{document}

